% Lösungen Einheit 4 von 4 – Ohne Zurücklegen (zusammenbau v0.1)
\einheitenkopf{Einheit 4 von 4 · Ohne Zurücklegen}

%% TODO Lösung der Erklärzeile 21g) fehlt in der Bank
\erg{21}{a) 8 Kekse, davon 3 mit Schoko \quad b) $\frac{3}{8} \cdot \frac{2}{7} = \frac{6}{56} = \frac{3}{28}$ \quad c) $\frac{5}{9} \cdot \frac{4}{8} = \frac{20}{72} = \frac{5}{18}$ \quad d) $\frac{2}{6} \cdot \frac{4}{5} = \frac{8}{30} = \frac{4}{15}$ \quad e) $\frac{5}{8} \cdot \frac{3}{7} = \frac{15}{56}$ \quad f) $\frac{4}{9} \cdot \frac{3}{8} = \frac{12}{72} = \frac{1}{6}$ \quad g) \ldots \quad h) erste Stufe grün $\frac{3}{7}$; nach rot: rot $\frac{3}{6}$; nach grün: rot $\frac{4}{6}$, grün $\frac{2}{6}$ \quad i) $\frac{3}{5} \cdot \frac{2}{4} + \frac{2}{5} \cdot \frac{1}{4} = \frac{6}{20} + \frac{2}{20} = \frac{8}{20} = \frac{2}{5}$ \quad j) $\frac{6}{7} \cdot \frac{1}{6} = \frac{1}{7}$ \quad k) $\frac{4}{9} \cdot \frac{3}{8} \cdot \frac{2}{7} = \frac{24}{504} = \frac{1}{21}$ \quad l) erste Stufe ohne Nuss $\frac{10}{15}$; nach ohne, Nuss: Nuss $\frac{4}{13}$ \quad m) $\frac{1}{9} + \frac{8}{9} \cdot \frac{1}{8} + \frac{8}{9} \cdot \frac{7}{8} \cdot \frac{1}{7} = \frac{3}{9} = \frac{1}{3}$ \quad n) mit: $\frac{3}{6} \cdot \frac{3}{6} = 0{,}25$; ohne: $\frac{3}{6} \cdot \frac{2}{5} = 0{,}2$ – ohne Zurücklegen ist sie kleiner. \quad o) Drei rote: $\frac{5}{15} \cdot \frac{4}{14} \cdot \frac{3}{13} = \frac{60}{2\,730}$, drei grüne: $\frac{3}{15} \cdot \frac{2}{14} \cdot \frac{1}{13} = \frac{6}{2\,730}$ – zehnmal so wahrscheinlich, nicht doppelt. \quad p) (Vanille; Pflaume), (Pflaume; Vanille), (Vanille; Vanille); $\frac{15}{18} \cdot \frac{14}{17} = \frac{210}{306} = \frac{35}{51} \approx 0{,}686$; Ereignis: beide Berliner haben die gleiche Füllung.}

21h)

\baumzwei{rot/\frac{4}{7},grün/\frac{3}{7}}{rot/\frac{3}{6},grün/\frac{3}{6}}{rot/\frac{4}{6},grün/\frac{2}{6}}


21l)

\baumdrei{Nuss/\frac{5}{15},ohne/\frac{10}{15}}{Nuss/\frac{4}{14},ohne/\frac{10}{14}}{Nuss/\frac{5}{14},ohne/\frac{9}{14}}{Nuss/\frac{3}{13},ohne/\frac{10}{13}}{Nuss/\frac{4}{13},ohne/\frac{9}{13}}{Nuss/\frac{4}{13},ohne/\frac{9}{13}}{Nuss/\frac{5}{13},ohne/\frac{8}{13}}

\erg{22}{a) Fehler: Beim dritten Zug ist der Nenner nicht weitergezählt. Richtig: $\frac{6}{9} \cdot \frac{5}{8} \cdot \frac{4}{7} = \frac{120}{504} = \frac{5}{21}$. \quad b) Die zuerst gezogene Kugel fehlt jetzt im Beutel; es gibt ein mögliches Ergebnis weniger. \quad c) erst eine rote, dann eine blaue Kugel \quad d) $\frac{48}{50} \cdot \frac{47}{49} \approx 0{,}921$}
